\section{Theory and Design of Experiment}
\label{sec:theory}

% PAGE TARGET: ~12 pages.
% Day 3 (5/27). HEAVIEST DAY.
% This is the CUED-style replacement for "Methodology" -- focus on the
% theoretical assumptions and how they translate to the design of our
% experiments. Heavy algebra and full derivations go to Appendix C.

\subsection{Physical Model of Laid Lines}
\label{sec:theory:physical}
% - Periodic thickness modulation from brass wires of the papermaking mould.
% - Effect on transmitted intensity (Beer-Lambert) vs reflectance (Lambertian
%   with orientation modulation).
% - Why DRP carries directional information that flat reflectance does not.

\subsection{Directional Reflectance Profile}
\label{sec:theory:drp}
% - Definition: $I(\theta, \phi)$ stack per pixel.
% - Sampling lattice: $N_\theta \times N_\phi$.
% - Background subtraction model.

\subsection{Stage 1 -- Per-Pixel Direction Estimation}
\label{sec:theory:direction}
% - Per pixel, fit dominant azimuth of the directional reflectance lobe.
% - Output: $\hat{a}(x,y) \in [-90,90)$ degrees mod 180.
% - Algorithm: aggregation over $(\theta,\phi)$ then trigonometric collapse
%   to azimuth (full equations in Appendix C).

\subsection{Stage 2 -- Trigonometric Patchwise Enhancement}
\label{sec:theory:trigmask}
% - Aggregate per-pixel azimuth over $P\times P$ patches with stride $S$
%   (default $512 \times 512$, stride 256).
% - Must respect 180-periodicity: average $(\cos 2A, \sin 2A)$, not $A$.
% - Output: enhanced grayscale where laid-line orientation locally dominates.

\subsection{Stage 3 -- Laid-Line Detection}
\label{sec:theory:detect}

\subsubsection{Gabor-based estimator}
% - Gabor filter bank over period set $\{6,8,\dots,40\}$ px.
% - Patchwise dominant period -> aggregated distribution.

\subsubsection{Multi-phi detector}
% - Bands across folio at multiple azimuths; per-band peak-find.

\subsubsection{Automatic line-direction detection}
% - Spectral power across rotated projections (NOT autocorrelation,
%   see commit caecc39).
% - Dual-band cross-validation rejects orientations with peak in only one band.
% - Polarity flag wire\_is\_darker for reflective vs transmitted-light images.

\subsection{Stage 5 -- Evaluation Metrics}
\label{sec:theory:metrics}
% Five outputs, each a JSON + companion PNG:
%
% (a) fit\_quality.json -- $R^2$ of harmonic model up to $k=4$.
% (b) interval\_distribution.json -- empirical peak-to-peak interval distribution.
% (c) self\_contrast.json -- $z = (\mathrm{score}-\mu_{\mathrm{null}})/\sigma_{\mathrm{null}}$
%     against rotated-image null distribution.
% (d) split\_half\_stability.json -- recompute on each half, compare.
% (e) wire\_width\_stats.json -- per-segment FWHM, median and CI.
%
% Full equations and derivations -> Appendix C.

\subsection{Design of Validation Experiments}
\label{sec:theory:design}
% Why we use:
% - Synthetic phantom sweeps (period, SNR, angle, sigma): isolate single
%   failure modes.
% - 9-folio real benchmark (4 manuscript stocks, varied conditions):
%   external validity.
% - Manual GT on two folios: collapse spreadsheet-folio mismatch ambiguity.
